\section{Power dynamics and the sporadic separation theorem}
\label{sec:power-dynamics}

The forward operators retain a finite dynamical invariant that is different
from the adjoint-closed algebra.  Let \(\tau:X\to X\) be a map of a finite
set and let \(P_\tau\) act on \(\Q^X\) by pullback,
\((P_\tau f)(x)=f(\tau x)\).  Write \(c_\ell(\tau)\) for the number of
directed cycles of length \(\ell\) in the functional graph of \(\tau\).

\begin{theorem}[Finite-map zeta formula]
\label{thm:finite-map-zeta}
For every finite map \(\tau:X\to X\),
\begin{equation}
 Z_\tau(z):=\det(I-zP_\tau)^{-1}
 =\prod_{\ell\geq1}(1-z^\ell)^{-c_\ell(\tau)}.
 \label{eq:finite-map-zeta}
\end{equation}
Thus \(Z_\tau\) depends only on the periodic cycles.  The rooted trees
feeding those cycles contribute no factor.
\end{theorem}

\begin{proof}
Order the vertices of each connected functional-graph component by
decreasing distance from its unique directed cycle, followed by the cycle
vertices.  In this basis \(P_\tau\) is block triangular.  Its off-cycle
diagonal block is nilpotent, while a cycle of length \(\ell\) contributes
the permutation matrix of an \(\ell\)-cycle.  Hence
\[
 \det(I-zP_\tau)
 =\prod_{\text{cycles }C}\det(I-zP_C)
 =\prod_{\text{cycles }C}(1-z^{|C|}).
\]
\end{proof}

For a finite group \(G\), put \(N_G=\exp(G)\) and
\[
 \mathcal N_G=\{p^a:p\text{ prime},\ a\geq1,\ p^a\mid N_G\}.
\]
Let \(\sigma_n:\Conj(G)\to\Conj(G)\) be \(C\mapsto C^n\).  The
\emph{ramified dynamical packet} is the unlabelled multiset
\begin{equation}
 \mathfrak Z_{\rm ram}(G)
 =\bigl\{\!\bigl\{Z_{\sigma_n}(z):n\in\mathcal N_G\bigr\}\!\bigr\}.
 \label{eq:ramified-dynamical-packet}
\end{equation}
It forgets both the prime-power labels and the ordering of the conjugacy
classes.  It is therefore weaker than the fully marked eidetic object.

\begin{theorem}[Sporadic dynamical separation]
\label{thm:sporadic-dynamical-separation}
For the twenty-six sporadic simple groups:
\begin{enumerate}
\item the map \(G\mapsto\mathfrak Z_{\rm ram}(G)\) is injective;
\item the cycle lengths occurring among all
      \(\sigma_n\), \(n\in\mathcal N_G\), lie in
      \(\{1,2,3,5\}\);
\item for each \(G\), the set of occurring cycle lengths equals
      \(\{1\}\) together with the orders of the good-label characters in
      \cref{thm:good-label-census};
\item a cycle of length \(3\) occurs exactly for
      \[
      J_1,\ J_3,\ J_4,\ Ly,\ O'N,\ Ru,
      \]
      and a cycle of length \(5\) occurs exactly for \(Ly\).
\end{enumerate}
\end{theorem}

\begin{proof}
The class power maps in the fixed CTblLib ordering give \(258\)
prime-power directions.  For each direction, cycle enumeration followed by
\cref{thm:finite-map-zeta} gives its exponent vector
\((c_1,c_2,c_3,c_5)\).  Sorting these vectors within each group produces
twenty-six distinct multisets, proving (1).  Direct enumeration gives
(2).  Comparing the cycle-order support with the independently reconstructed
good-label characters of \cref{thm:good-label-census} proves (3).
The same exact enumeration gives the two lists in (4).
\end{proof}

\begin{corollary}[Odd cyclotomic detection]
\label{cor:dynamics-cyclotomic-detection}
Among the sporadic simple groups, the ramified dynamical packet detects
every nonrational center field in \(\Pram(G)\).  Cubic cycles occur exactly
for the six groups carrying a \(\Q(\zeta_3)\)-factor, while the unique
quintic cycle occurs for the unique group carrying a
\(\Q(\zeta_5)\)-factor.
\end{corollary}

\begin{proof}
Combine \cref{thm:sporadic-dynamical-separation} with
\cref{thm:sporadic-census,thm:good-label-census}.
\end{proof}

\begin{remark}
The equality of supports in
\cref{thm:sporadic-dynamical-separation}(3) is a census theorem, not a
general equivalence for finite groups.  The general implication supplied
by \cref{lem:orbit-exponent-bound} is one-sided: character order is bounded
by permutation order on block support.
\end{remark}
